\section{全重みで増大する厚みと明示的な量子代数}
\label{sec:allm-notes-v33}

この節の結果は公表稿IX\cite{PaperIX}に証明をまとめる。
一般重みの厚みは仮説ではなく、全ての整数 $m\ge3$ に対する定理となった。
今後は、この機構と明示的な代数を使って別の変形問題を判断する。

\subsection{全方程式と、重み三だけに残る慣性成分}
\[
 X_m=\mathbb P_{\mathrm{stack}}(1,1,1,m),\quad
 S_d=\mathbb C[x,y,z]_d,\quad N_m=\binom{m+2}{2}
\]
と置く。$m\ge3$ では標準形の座標は
\[
 c\in S_3,\qquad D\in\mathfrak{sl}_3,\qquad
 [F]\in S_{m+1}^{\oplus3}/\mathbf xS_m
\]
であり、全Jacobi式は
\[
 D(c)=0,\qquad
 \frac{m+3}{3}F(c)-\frac m3(\mathbf x\times D)\cdot F-c\operatorname{div}F=0.
\]
変数数と独立な二次式の数は
\[
 d_m=m^2+6m+26,\qquad q_m=10+\binom{m+5}{2}.
\]
二次式の個数を余次元と取り違えない。
$m=1,2$ では二次場の係数に追加の $w$ 項があるので、上の標準形をそのまま使わない。

標準 $(m+3)$ 項の代数 $\Lambda_m$ は
$\dim\Lambda_m=\binom{m+6}{4}+15$ であり、そのminiversal基底は
このJacobiスキームの原点完備化である。
証明に再利用する機構は、Euler発散による一意な持ち上げ
\[
 L(U)=U-\frac1{m+3}E\wedge\delta U
\]
と、全Jacobi軌跡上の量子化、二次生成性に対するNakayamaの補題である。
全高次障害を個別に計算する必要はない。

$HH^2(\Lambda_m)$ は $d_m$ 次元である。
$m\ge4$ では $HH^3(\Lambda_m)$ は $q_m$ 次元で、一次方向間の括弧の像が全体を張る。
一方 $m=3$ では $\dim HH^3=41$、括弧の像は三十八次元である。
この差には、重み六の $w^2$ と、二つの非自明慣性成分が寄与する。
一般重みで慣性成分を数える際は、各点 $B\mu_m$ 上の法束行列式の文字
$\chi^3$ に対して\emph{全安定化群}の不変部分を取る。
個々の慣性元の位数だけを見ると、$m$ が三の倍数の場合に誤った寄与を加えてしまう。

\subsection{一般点のArtin環を決める三つの情報}
基準点は全重みを通じて
\[
 c_0=xy(x+y),\qquad D_0=(2x+5y)\partial_z,\qquad F=0
\]
に取れる。この点を含む適切な開集合で、全完備局所環は
\[
 \boxed{\mathbb C[[s_1,\ldots,s_{N_m+8},a,b,u]]/(ab,b^2,a^{m+2}u)}.
\]
その証明を再利用するには、次の三点を保つ。
第一に、基底 $D(c)=0$ の切片は $m$ によらず $(ab,b^2)$ であり、
平方零方向 $b$ を消去せずに混合式へ渡す。
第二に、混合核 $N_m+1$ 方向のうち、次数 $m$ の $w$ 移動の
$N_m$ 方向は非被約基底上でも実際の核の自由な直和因子である。
残る一方向は $F=(0,0,x^{m+1})$ で、$z$ 重量は $-1$。
余核の四方向の重量は $m+1,m+2,m+2,m+3$ である。
従って $a$ の障害は $m+2$ 次以上、$b$ の障害は零となる。

第三に、$b=0$ の三角形 $c=xy(x+y+az)$ の三交点での評価を使う。
Schur消去後の出力を
\[
 r_0x^2z^{m+1}+r_1xz^{m+2}+r_2yz^{m+2}+r_3z^{m+3}
\]
とすると、交点評価から $r_1=ar_0,r_2=r_3=0$ となる。
消去行列の重量から $r_0$ の $a$ 次数は $m+2$ 以下であり、
$a\ne0$ 側の階数が正確に一つ増えるので、$a^{m+2}$ 係数は必ず非零である。
ここが有限個の重みの検算を全重みの証明へ置き換えた箇所である。

引戻し成分の一般点では $u$ が単元なので
\[
 K[a,b]/(a^{m+2},ab,b^2),\qquad
 1,a,\ldots,a^{m+1},b
\]
が得られる。成分次元は $N_m+9$、長さは $m+3$、socle次元は二である。
被約化は滑らかであっても、この厚みは一般点に残る。

\begin{center}\small
\begin{tabular}{rrrrr}\toprule
$m$&$d_m$&$q_m$&引戻し成分次元&一般点の長さ\\\midrule
3&53&38&19&6\\
4&66&46&24&7\\
5&81&55&30&8\\
6&98&65&37&9\\\bottomrule
\end{tabular}
\end{center}

primary分解と交叉の厚みも全重みで一様である。
\[
 (ab,b^2,a^{m+2}u)=(a^{m+2},ab,b^2)\cap(b,u).
\]
二つの被約枝は引戻し型と $L(1,1,1,m)$ 型であり、被約交叉は滑らか。
primary成分どうしの交叉は長さ $m+2$、一般超平面切断は長さ $m+4$ である。
局所環はCohen--Macaulayだが完全交叉ではない。

Poissonスキームは斉次錐なので、原点での形式化後にも対応する成分と一般点の長さが残る。
証明では、次数付き整域の完備化も最小斉次項の積により整域となること、
原点の局所環からの忠実平坦性、組成列の基底変換を用いる。
従って $\Lambda_m$ のminiversal基底にも、一般点の長さが任意に大きい成分が現れる。
この引戻し成分の原点基底の計算だけでは、固定した非零量子代数の変形空間は決まらない。
接線共鳴の固定量子化については、公表稿Xで別途形式弧移送により決定した。

\subsection{全重みで使える明示的AS正則4-Calabi--Yau代数}
二変数の $c_3,L_1,f_{m+1}$ を任意に取り、
\[
 \begin{split}
 \delta_m(x)&=c_y/3+mxL/(3(m+3)),\\
 \delta_m(y)&=-c_x/3+myL/(3(m+3)),\\
 \delta_m(w)&=f-mLw/(m+3)
 \end{split}
\]
と置く。$\operatorname{div}\delta_m=0$ であり、
\[
 S_m=\mathbb C[x,y,w][z;\hbar\delta_m]
\]
は引戻し族の適合座標表示を量子化する。
PBW基底は $x^iy^jw^kz^\ell$、Hilbert級数は $((1-t)^3(1-t^m))^{-1}$。
noetherian整域で、AS正則次元四、通常の4-Calabi--Yau、Gorenstein parameter $m+3$ である。
星積は
\[
 a\star b=\sum_{r\ge0}\frac{\hbar^r}{r!}
                 (\partial_z^r a)(\delta_m^r b)
\]
であり、各多項式対について有限和になる。
抽象的な形式存在を超えて、任意の複素数 $\hbar$ で座標環を持つ点が有用である。

標準列 $(\mathcal O,\ldots,\mathcal O(m+2))$ はfull strongであり、
$B_m=(S_m)_{(0,1,m,m+1)}$ は階数 $2m^2+6m+20$ の四対象代数を与える。
$\operatorname{Perf}(B_m)$ は $D^b(\operatorname{qgr}S_m)$ の許容部分圏である。
$m=3$ は元のVK問題に対する、新しい具体的な標準六項実現である。
$m>3$ で $B_m$ 自身の変形が無障害であることや、外側との接空間同型は未確認である。
この二つは今回の量子化の存在からは従わない。

\subsection{接線共鳴は全重みの成分系列を作る}

$q=xz-y^2$、$N_0=x\partial_y+2y\partial_z$ とする。
$0\le j\le\lfloor(m-1)/2\rfloor$、$r=m-2j>0$ に対して
\[
 c=xq,\qquad D=a_jN_0,\quad a_j=-\frac{4m+3-6j}{2m},\qquad
 F_j=q^j(0,yz^r/2,z^{r+1})
\]
を基準点とする。$(c,D)$ の七次元軌道上の混合核束の閉包は、
次元 $N_m+8$ の既約成分 $\mathscr E_{m,j}$ である。
全 $m,j$ について一般点の局所環と $F=0$ の全形式近傍が
\[
 K[\gamma]/(\gamma^{j+1}),\qquad
 \mathbb C[[s_1,\ldots,s_{N_m+7},\kappa,\gamma,u]]/
                    (\kappa u,\gamma^{j+1}u)
\]
と決まった。固定した $m$ で相異なる $\lceil m/2\rceil$ 成分を与え、
長さは $1,\ldots,\lceil m/2\rceil$ を尽くす。
$j=0$ は滑らか、$j\ge1$ は一般点でも非被約である。
$u=0$ の枝は $N_m+9$ 次元の $L(1,2,m)$ 成分となる。

再利用すべき証明の機構は、次の具体的な四点である。
基底 $D(c)=0$ は七軌道方向を除くと
\[
 c=x(xz-y^2+\gamma z^2),\qquad
 D=a((x+2\gamma z)\partial_y+2y\partial_z)
\]
の二変数となる。混合核を $F_x=0$ のEuler代表で書くと
\[
 (F_y,F_z)=(xA+yB,2yA+2zB),\qquad
 N_0(A)=(T-2j)B,\quad T=2x\partial_x+y\partial_y
\]
である。実際の $w$ 移動は
\[
 A=-\tfrac12(T+1-2j)g,\qquad B=-\tfrac12N_0(g)
\]
となり、階数 $N_m$ の自由な直和因子をなす。
例外的な $T$ 重量 $2j-1\to2j$ の二重対角行列は単射で、
余核の一元が $q^jz^r$ である。

混合余核の $H_0=x\partial_x-z\partial_z$ 重量は
$-m-3,-m-2,-m-1,1-r$ である。
重量を保つ単元ブロックの逆行列は $\gamma$ の有限多項式なので、
残余列の形は正確に
\[
 f(a),\qquad \gamma^{j+1}g(a),\qquad0,\qquad\gamma^{j+2}h(a)
\]
へ制約される。$f$ は $a-a_j$ の単元倍である。
二つの特異点 $[0,\pm\sqrt\gamma,1]$ での評価から
$h(a_j)=-g(a_j)$、横断二次曲線の場合の混合核を次数降下で全て解くことから
$g(a_j)\ne0$ と分かる。この二つを合わせて全理想が
$(a-a_j,\gamma^{j+1})$ と確定する。
この方法は最初の障害次数だけでなく、高次の隠れた条件がないことまで保証する。

primary分解は
\[
 (\kappa u,\gamma^{j+1}u)=(u)\cap(\kappa,\gamma^{j+1})
\]
である。被約交叉は滑らかな $N_m+7$ 次元、primary成分の交叉には
長さ $j+1$ の厚みが残る。
$r=0$ は $F_x$ の消去が成立しないため、この成分定理の対象から外す。
量子化と中心の定理は、この端点も含めて別に成立する。

\subsection{重み三の第二共鳴と、二次で止まる方向}

$m=3,j=1$ の同じ成分は、$qz$ 型と等変同値な $y^2z$ 型で
\[
 D=-\tfrac32N_0,\qquad F=(0,y^3z/2,y^2z^2)
\]
とも書ける。この点の完備局所環は
$\mathbb C[[s_1,\ldots,s_{18},b]]/(b^2)$ である。
直接の有理係数消去でも、この全局所環を独立に確認した。
特に
\[
 \begin{split}
 c_b&=x(xz-y^2-bz^2/3),\\
 D_b&=(0,-3x/2+bz,-3y),\\
 F_b&=(0,y^3z/2-byz^3/12,y^2z^2+bz^4/6)
 \end{split}
\]
について $D_b(c_b)=0$、$M(c_b,D_b;F_b)=-7b^2xz^5/36$ である。
接触を解く一次方向が二次で止まることを、今後の障害計算の最小例として使える。
この成分を加え、重み三の認定済み大域成分は少なくとも十個となった。

非零の接線パラメータの四共鳴は、滑らかな十七変数を除いて
\[
 \begin{array}{c|c}
 a=2&(kv_1,\ldots,kv_{10})\\
 a=-1&(kv_1,\ldots,kv_4)\\
 a=-5/2&(kv,bv)\\
 a=-3/2&(kv,b^2v)
 \end{array}
\]
という局所方程式を持つ。最初の二つは $R(1,5),R(2,4)$ との交叉である。
非共鳴で混合核の次元が十であっても、常に全てを $w$ 移動と同一視してはいけない。
$a=1$ では $p=(q,0,0)$ となり、移動核には $x^3$ が入る。
欠ける方向 $x^3N_0$ は
$p_a\times\nabla(x^3)=(1-a)x^3N_0$ の有限極限として得られ、
全空間はこの点でも滑らかな十九次元である。

\subsection{非線形CGG族のNakayama自己同型は全て決まる}

$m=2j+r\ge3$、$j,r\ge0$ とし、$P=q^jz^r$ または $y^{2j}z^r$ とする。
\[
 N=N_0-P\partial_w,\qquad d=m+r+1,
\]
\[
 X_\alpha=\alpha x\partial_x+(\alpha+1)y\partial_y+(\alpha+2)z\partial_z
                   +(m\alpha+d)w\partial_w
\]
に対する有限CGG積を考える。
$T=1+sN$ は星積の自己同型で、変数を $x,y,z,w$ の順に加えると反復Ore表示を得る。
全 $\alpha,s$ でnoetherian整域かつAS正則四次元であり、
Hilbert級数は $((1-t)^3(1-t^m))^{-1}$、Gorenstein parameterは $m+3$ となる。
そのNakayama自己同型は全次数で
\[
 \boxed{\mu=T^\nu,\qquad \nu=(m+3)\alpha+m+r+3=\operatorname{div}X_\alpha-1}
\]
と決まる。従って $s\ne0$ で通常の4-Calabi--Yauとなるのは
$\alpha=-(m+r+3)/(m+3)$ の場合に限る。

再利用できる剛性は、次数一で生成される三変数部分代数を各元ごとに固定する
次数付き自己同型が恒等になることである。
$\psi(w)=uw+b$ を、$N$ 不変な奇数次数の $J$ に作用させる。
$x,q$ との可換性から $N(\psi J-uJ)=0$、$y$ との交換関係から
その $H_0$ 重量が零となる。
奇数次数・零重量の空間上の $N_0$ は二重対角の正則行列であり、従って差は零である。
$J\bmod x$ が非零であることから $u=1,b=0$ となる。
この議論は $\alpha=0,-1$ も含み、Ore定理だけでは未決の最後の平行移動項を消す。

Lecoutre--Sierraの三変数Nakayama公式を用いる部分と、この四変数剛性を区別する。
通常のCY性を発散零の一次テンソルだけから推論したわけではない。
$s\ne0$ の次数付き自己同型群は
$\mathbb C^\times\times(\mathbb C,+)$ であり、作用は
重み付きスカラー変換と $T^t$ である。
全ての族は標準 $(m+3)$ 項のfull strong列と四対象の許容成分を持つ。

\subsection{多項式中心の三分岐とDu Val曲面}

$q$ 型で
\[
 c_r=(-1)^r\frac{(2r)!!}{(2r+1)!!},\qquad
 J=x^{r+1}w+q^j\sum_{k=0}^r\binom rk\frac{q^{r-k}y^{2k+1}}{2k+1}
\]
とすると、
\[
 \Delta=(J^2+c_r^2q^{m+r+1})/x^{r+1}
\]
は多項式で、
\[
 \ker N=\mathbb C[x,q,J,\Delta]
 \simeq\mathbb C[X,Q,U,V]/(X^{r+1}V-U^2-c_r^2Q^{m+r+1})
\]
である。$x$ を反転した核だけでは生成性は決まらない。
$x=0$ でのcusp表示と $\Delta$ の非零の $w$ 一次係数から
$B/xB\to A/xA$ の単射性を示し、$x$ に関する飽和で核全体を得るのが重要である。

Calabi--Yauのパラメータでは、Poisson中心と量子中心が同じ部分環として
$\ker N\cap\ker X_\alpha$ に一致する。
$2(m+3)H=- (m+3)X_\alpha$ の重量は
\[
 x:m+r+3,\quad q:2r,\quad J:r(m+r+1),\quad
 \Delta:(r-1)m+(r-3)(r+1)
\]
なので、$r\ge2$ では中心は定数環、$r=0$ では $\mathbb C[q,J]$ である。
$r=1$ では
\[
 S=q^2\Delta,\quad P=2xJ\Delta^{(m+3)/2},\quad
 R=q(x^2\Delta+J^2)\Delta^{(m+3)/2}
\]
が中心全体を生成し、その唯一の関係は
\[
 \boxed{R^2=SP^2+\frac{16}{81}S^{m+3}}
\]
である。これは $D_{m+4}$ 型のDu Val曲面である。
生成性は、$\mathbb C[x,q,\Delta]\oplus J\mathbb C[x,q,\Delta]$ の
零重量単項式を指数の偶奇で完全に分けることで証明する。

$y$ 型との等変変換も多項式で与えられる。
\[
 \lambda_{j,r}=(-1)^j\frac{(2j-1)!!(2r+1)!!}{(2j+2r+1)!!},\qquad
 w'=\lambda_{j,r}^{-1}(w+b_{j,r}),
\]
\[
 b_{j,r}=(J_0^{(y)}-\lambda_{j,r}J_0^{(q)})/x^{r+1}\in S_m.
\]
多項式性は二つのTaylor表示の $r$ 次までの係数一致から従う。
特に重み三の第二共鳴は $w'=-5w-yz^2$ で $qz$ 型となり、その量子中心は $D_7$ 型である。

一般点の長さと中心の型が同時に読める系列は
\[
 m=2\ell-1,\quad j=\ell-1,\quad r=1:\qquad
 \text{成分の長さ }\ell,\quad \text{量子中心 }D_{2\ell+3}\quad(\ell\ge2)
\]
である。前者は可換スタック側のPoissonスキームの一般点長さである。
公表稿Xでは形式弧移送と軌道方向の除去により、固定した非零量子化の
標準有限代数のminiversal環も正確に $\mathbb C[[\gamma]]/(\gamma^\ell)$ と同定した。

\subsection{共通の関数斜体と、同じものにしない不変量}

全 $j,r,\alpha$ について、$s\ne0$ なら次数零の関数斜体は
$\mathbb C(C)$ 上の第一Weyl斜体である。
\[
 t=x^{-1}\star y,\quad u=q/x^2,\quad v=J/x^d,\quad
 [t,u]=2su,\quad[t,v]=dsv,\quad[u,v]=0
\]
とし、$d=2h+1$ に対し $C=v^2/u^d$、$\eta=v/u^h$ と取れば
$\partial=\eta^{-1}t/s$ は $[\partial,\eta]=1$ を満たす。
有理Poisson中心と量子関数斜体の中心はともに $\mathbb C(C)$ である。
したがって多項式中心が定数・平面・Du Val曲面のどれであるかは、
この射影的な関数斜体の同型型には残らない。
異なる $m$ の射影圏は $K_0$ の階数 $m+3$ で区別できる。
同じ $m$ の異なる $j$ の射影圏を座標環の中心だけで区別することはできない。
公表稿Xの固定変形環と反標準Hochschild--Serre環を用いると、本族内部の全ての $j$ を区別できる。

補助次数 $(0,1,2,d)$ による変数の拡大縮小は、非零の $s$ を互いに移す。
また $X_{\alpha+\beta}=X_\alpha+\beta E$ とVandermonde恒等式から
\[
 a\star_{\alpha+\beta}b=T^{\beta\deg b}(a)\star_\alpha b
\]
となる。固定した $j,r$ の全 $\alpha$ はZhang twistで結ばれ、射影圏は同値である。
通常のCY性は、この同じ射影圏の中の特定の次数付き代表を選ぶ条件でもある。

\subsection{$c=0$ と残る大域分類}

重み三で $c=0$ の被約部分は、ちょうど次元 $35,20,18$ の三成分を持ち、
全スキームではそれぞれ $\operatorname{PB}(\mathbb P^2)$、$L(1,1,4)$、
$L(1,1,1,3)$ に含まれる。
従って $c=0$ に支持される未知の大域成分はない。
ここでの十八次元部分は、接線型の二つの大域成分とは別である。
再利用する事実は、巡回型の $D$ に対して
$(\mathbf x\times D)\cdot F=0$ の全多項式解が $F=f\mathbf x+gD$ と書けることである。
商体上の解を高さ一の局所環で確認し、正規性で多項式に戻す。

接線型の非零 $D$ で残っていた第二共鳴は完了した。
大域分類の残りは非被約三次形式、その他の特異三次形式の特殊係数、
およびそれらの全混合核の大域成分への所属である。
十成分で分類が尽きることと最大次元は、なお未確定である。
